\documentclass{article}
% template.tex:
\usepackage{ifthen}
\usepackage[utf8]{inputenc}
\usepackage{graphics}
\usepackage{graphicx}
\usepackage{textcomp}
\usepackage{tcolorbox}
\usepackage{hyperref}
\usepackage{svg}
\usepackage{multicol}
\usepackage[section]{placeins}
\usepackage[backend=biber,style=numeric]{biblatex}
\addbibresource{sources.bib}

\pagestyle{plain}
\voffset -5mm
\oddsidemargin  0mm
\evensidemargin -11mm
\marginparwidth 2cm
\marginparsep 0pt
\topmargin 0mm
\headheight 0pt
\headsep 0pt
\topskip 0pt
\textheight 255mm
\textwidth 165mm

\newcommand{\duedate} {}
\newcommand{\setduedate}[1]{%
\renewcommand\duedate {Due date:~ #1}}
\newcommand\isassignment {false}
\newcommand{\setassignment}{\renewcommand\isassignment {true}}
\newcommand{\ifassignment}[1]{\ifthenelse{\boolean{\isassignment}}{#1}{}}
\newcommand{\ifnotassignment}[1]{\ifthenelse{\boolean{\isassignment}}{}{#1}}

\newcommand{\assignmentpolicy}{
\begin{tcolorbox}[colback=gray!10, colframe=black,
                  title={HPC Lab for CSE 2025 --- Submission Instructions},
                  fonttitle=\bfseries,
                  sharp corners,
                  boxrule=0.8pt,
                  left=5pt, right=5pt, top=5pt, bottom=5pt]
% \small
\textbf{Following these instructions is mandatory}: Submissions that do not
comply \textbf{will not be considered}.
\begin{itemize}
\item \textbf{Submission Platform}:
      Assignments must be submitted via
      \href{https://sam-up.math.ethz.ch/?lecture=401-3670-00}{SAM-UP}.
\item \textbf{Required Files}: Your submission must include:
      \begin{itemize}
      \item \textbf{Report}: A PDF summarizing your results and observations,
            based on this \LaTeX\ template.
            \begin{itemize}
            \item File name:
                  \texttt{project\_number\_lastname\_firstname.pdf}
            \item You are allowed to discuss all questions with anyone you like;
                  however: (i) your submission must list anyone you discussed
                  problems with and (ii) you must write up your submission
                  independently.
            \end{itemize}
      \item \textbf{Source Code Archive}: A \textbf{compressed tar archive},
            named \texttt{project\_number\_lastname\_firstname.tar.gz},
            containing:
            \begin{itemize}
            \item All source code files.
            \item Build files and scripts (e.g., \texttt{Makefile}, batch job
                  scripts, plot scripts, etc).
                  If you modified any provided files, ensure they still work as
                  expected.
            \item The full \LaTeX\ source of your report including all needed
                  files to build the report (figures, listings, etc).
            \end{itemize}
      \item \textbf{Size Limit}: The total submission (PDF + archive) must not
            exceed \textbf{25 MB}.
      \item \textbf{File Naming Rules}: Use only \texttt{ASCII} characters
            (avoid spaces, special symbols, and non-English letters).
      \end{itemize}
\item \textbf{Grading Process}: The TAs will evaluate your project based on:
      \begin{itemize}
      \item Your report write-up.
      \item Your code implementation.
      \item Benchmarking your code's performance.
      \end{itemize}
\end{itemize}
\end{tcolorbox}
}
\newcommand{\punkte}[1]{\hspace{1ex}\emph{\mdseries\hfill(#1~\ifcase#1{Points}\or{Points}\else{Points}\fi)}}


\newcommand\serieheader[6]{
\thispagestyle{empty}%
\begin{flushleft}
\includegraphics[width=0.4\textwidth]{ETHlogo_13}
\end{flushleft}
  \noindent%
  {\large\ignorespaces{\textbf{#1}}\hspace{\fill}\ignorespaces{ \textbf{#2}}}\\ \\%
  {\large\ignorespaces #3 \hspace{\fill}\ignorespaces #4}\\
  \noindent%
  \bigskip
  \hrule\par\bigskip\noindent%
  \bigskip {\ignorespaces {\Large{\textbf{#5}}}
  \hspace{\fill}\ignorespaces \large \ifthenelse{\boolean{\isassignment}}{\duedate}{#6}}
  \hrule\par\bigskip\noindent%  \linebreak
 }

\makeatletter
\def\enumerateMod{\ifnum \@enumdepth >3 \@toodeep\else
      \advance\@enumdepth \@ne
      \edef\@enumctr{enum\romannumeral\the\@enumdepth}\list
      {\csname label\@enumctr\endcsname}{\usecounter
        {\@enumctr}%%%? the following differs from "enumerate"
        \topsep0pt%
        \partopsep0pt%
        \itemsep0pt%
        \def\makelabel##1{\hss\llap{##1}}}\fi}
\let\endenumerateMod =\endlist
\makeatother


\usepackage{listings}
\usepackage{xcolor}

% Define colors for syntax highlighting
\definecolor{codegreen}{rgb}{0,0.6,0}
\definecolor{codegray}{rgb}{0.5,0.5,0.5}
\definecolor{codepurple}{rgb}{0.58,0,0.82}
\definecolor{backcolour}{rgb}{0.95,0.95,0.95}

% Configure C style
\lstdefinestyle{cstyle}{
    backgroundcolor=\color{backcolour},
    commentstyle=\color{codegreen},
    keywordstyle=\color{blue},
    numberstyle=\tiny\color{codegray},
    stringstyle=\color{codepurple},
    basicstyle=\ttfamily\footnotesize,
    breakatwhitespace=false,
    breaklines=true,
    captionpos=b,
    keepspaces=true,
    numbers=left,
    numbersep=5pt,
    showspaces=false,
    showstringspaces=false,
    showtabs=false,
    tabsize=2,
    language=C
}

\lstset{style=cstyle}


\newfloat{Listing}{htpb}{lop}
\DeclareCaptionSubType{Listing}


\begin{document}




 

   





 

\textbf{Randomness Tests} * \textbf{Uniformity (\(\chi^{2}\)):} Divides interval into B bins, compares observed (\(n_{i}\)) to expected (\(E_{i} = M/B\)) counts: \(\chi^{2} = \frac{1}{B}\sum_{i = 1}^{B}\frac{\left( n_{i} - E_{i} \right)^{2}}{E_{i}}\). Small \(\chi^{2}\) indicates agreement. * \textbf{Correlation:} Measures time correlation \(C(k) = \frac{\langle x_{i}x_{i + k}\rangle_{i} - \langle x_{i}\rangle_{i}\langle x_{i + k}\rangle_{i}}{\langle x_{i}x_{i}\rangle_{i} - \langle x_{i}\rangle_{i}\langle x_{i}\rangle_{i}}\). \(C(k) = 0\) for no correlation. \(\langle x_{i}x_{i + k}\rangle_{i} = \frac{1}{B - k}\sum_{i = 1}^{B - k}x_{i}x_{i + k}\)

\subsection{Equations of Motion}\label{equations-of-motion}    \subsubsection{Lagrangian formulation}\label{lagrangian-formulation}  

\subsubsection{Hamiltonian formulation}\label{hamiltonian-formulation}  

\subsection{Schrödinger Equation}\label{schrödinger-equation}  

\subsection{Thermodynamics}\label{thermodynamics}    





 

 

 

 

\section{Molecular Dynamics}\label{molecular-dynamics}   

\subsubsection{Degrees of Freedom}\label{degrees-of-freedom}    

\subsubsection{Classical Atomistic Force fields}\label{classical-atomistic-force-fields}   

     

\textbf{Integrators} * Integrate \(\frac{d^{2}r_{i}(t)}{dt^{2}} = \frac{f_{i}}{m_{i}}\). * \textit{Leapfrog}: \(v\left( t_{n} + \frac{\Delta t}{2} \right) = v\left( t_{n} - \frac{\Delta t}{2} \right) + \frac{f\left( x_{n} \right)}{m_{i}}\Delta t + \mathcal{O}\left( \Delta t^{3} \right)\), \(x\left( t_{n} + \Delta t \right) = x\left( t_{n} \right) + v\left( t_{n} + \frac{\Delta t}{2} \right)\Delta t + \mathcal{O}\left( \Delta t^{3} \right)\). Timestep: usually order of \(10^{- 15}\) * \textit{Euler, Runge-Kutta}. Balance performance/accuracy. Low order (leapfrog/verlet) lead to noisy forces. Algorithms without velocities cannot couple to thermostats.

\textbf{Metropolis Monte Carlo} * Samples configurations with canonical (Boltzmann) weight. Trial move accepted with probability: \(P = \min(1,e^{- \beta\Delta V})\), where \(\Delta V\) is the (pot.) energy difference. \(\beta = \frac{1}{k_{B}T}\). For uphill moves \((\Delta V > 0)\), accept if \(s < e^{- \beta\Delta V}\) with \(s \in U\lbrack 0,1\rbrack\).

\textbf{Bond constraints} * Timestep limited by fastest motion (~5fs bond vibration). Rule of thumb \(\Delta t = \tau/10\). * Constraining bonds to fixed length increases timestep (~2fs). Uninteresting for larger systems. * Algorithms: \textit{SHAKE, RATTLE, SETTLE, LINCS}. SHAKE is iterative.

\subsubsection{Boundary Conditions}\label{boundary-conditions}  \textbf{Spatial Boundary Conditions} * Finite size effects (surface). Worse for microscopic sim. * Vacuum effects: Surface tension, larger charge interactions. Remedy: Implicit/explicit solvents, reduce charges. * Boundary layer: Restrict motion (harmonic potential), add ext. force.

\textbf{Periodic Boundary Conditions (PBC)} * Avoids surface effects. Simulates infinite lattice. Particle exiting enters opposing face. * Interactions evaluated using lattice-sum methods (Fourier series). * Issue: Self-interaction if period too small (\(R_{c} < L/2\)). Large computation effort. Truncate non-bonded interactions at \(R_{c}\). * Easiest shape: boxes. Other possible.

\textbf{Thermodynamic Boundary Conditions} * Fix state variables. Imposed by ensembles. * Soft coupling desired (fixed T long-term, but fluctuates). Negative-feedback mechanism.

\textbf{Instantaneous pressure} * From \textit{virial theorem}: \(P = \frac{2(K–W)}{3\mathcal{V}}\), \(W = \frac{1}{2}\sum_{i = 1}^{N}r_{i} \cdot F_{i}\). * Barostating by isotropic coordinate scaling: Scale edges if P differs from target. * \(P = \frac{Nk_{B}T}{\mathcal{V}} - \langle\frac{2W}{3\mathcal{V}}\rangle\). Sum W using minimum-image vector.

\subsubsection{Thermostating, Barostating}\label{thermostating-barostating}  

\textbf{Temperature constraining} * Modify Newtons eqns: \(\dot{p} = - \nabla_{r}V–\zeta_{T}p\), \(\zeta_{T} = \frac{- \nabla_{r}V \cdot \dot{p} - F \cdot \dot{r}}{2K}\). Hard constraint on K (T). * Isokinetic ensemble for momenta. * Implementation (\textit{velocity scaling}): * \textbf{Hoover-Evans}: scale \(v\left( t + \frac{\Delta}{2} \right)\) by \(\lambda = \left( \frac{K\left( t - \frac{\Delta}{2} \right)}{K\left( t + \frac{\Delta}{2} \right)} \right)^{\frac{1}{2}}\). \(\dot{K} = 0\). \textbf{Woodcock}: scale \(v\left( t + \frac{\Delta}{2} \right)\) by \(\lambda = \left( \frac{K_{\text{target}}}{K\left( t + \frac{\Delta}{2} \right)} \right)^{\frac{1}{2}}\). Target K maintained, no drift.

\textbf{Berendsen Thermostat (weak coupling)}: * Not hard constraint, exponential relaxation of \(T \rightarrow T_{0}\). * \(\dot{T} = \frac{T_{0} - T}{\tau_{t}}\), \(\tau_{T}\): relaxation time. * \(T = 2\frac{K}{N_{\text{df }}k_{B}},K = \frac{1}{2}\sum m_{i}v_{i}^{2}\). * Velocity scaling each step: \(\lambda = \sqrt{1 + \frac{\Delta t}{\tau_{T}}\left( \frac{T_{0}}{T} - 1 \right)}\) * Effect: smooth \(T\) control, exp. relaxation, reduced fluctuations. * Ensemble ill-defined (depends on \(\tau_{T}\), \(C_{V}\)), ok for equilibration, not strict NVT 

 \textbf{Thermostating via stochastic methods} * \textit{Collision thermostat}: Randomly picks atoms, reassigns velocity from Maxwell-Boltzmann. * \textit{Langevin thermostats}: Add stochastic force (friction + virtual collisions). * Non-deterministic, perturbs dynamics. Needs good random numbers. Realistic exponential T relaxation, robust, efficient, homogeneous thermalization.

\textbf{Pressure constraining} * Modify eqns: \(\dot{r} = \nabla_{p}K + \chi r\), \(\dot{\mathcal{V}} = 3\chi\mathcal{V}\), \(\dot{p} = - \nabla_{r}V - \chi p\). * \(\chi\) from \(3P\mathcal{V} = 9P\chi\mathcal{V} = 2\sum_{i = 1}^{N}\frac{p_{i} \cdot {\dot{p}}_{i}}{m_{i}} - \sum_{i = 1}^{N}\left( F_{i} \cdot {\dot{r}}_{i} + r_{i} \cdot {\dot{F}}_{i} \right)\). Difficulty in virial derivative.



\textbf{Barostating via stochastic methods} * \textit{Langevin barostat}. Adds frict. to \(\mathcal{V}^{¨}\), stoch. V fluctua.

\textbf{Grand-canonical simulations} * N varies, \(\mu\) constant on average. Discrete variation. \(\mu = \frac{\partial F}{\partial N}\). * \textit{Van Gunsteren weak-coupling}: Estimate excess \(\mu\) (\(\mu^{ex} = - \beta^{- 1}\ln\langle e^{- \beta E_{test}}\rangle_{N}\)), adjust N: \(\frac{dN(t)}{dt} = \mathrm{\text{round}}\left( \alpha_{N}\left( e^{- \beta\mu^{ex}}–e^{- \beta\mu^{0}} \right) \right)\). Equilibration period after adjustment.

\textbf{Selection of coupling times} * Simultaneous barostating/thermostating: \(\tau_{T} < \tau_{P}\). Otherwise unphysical fluctuations. * P highly fluctuating (\(P \sim K - W\)).



\section{Stochastic Dynamics}\label{stochastic-dynamics}   \subsection{Calculation of Properties}\label{analysis-of-properties}  

\subsubsection{1. Preprocessing}\label{preprocessing}   

\subsubsection{2. Statistical analysis}\label{statistical-analysis}      

\subsubsection{3. Calculation of Properties}\label{types-of-properties}     

\subsubsection{4. Interpretation of results}\label{interpretation-of-results} 

\textbf{Ergodicity, equilibration} * Ergodic theorem holds for \(t \rightarrow \infty\). Not necessarily true for finite trajectories. * Discard initial period (equilibration) to account for non-representative start. * Equilibration possible in realistic time if initial config is reasonable (X-ray/NMR). * Improve avg estimate with multiple sims/initial conditions.

\textbf{Block Averaging}: \(\varepsilon = c\frac{\sigma}{\sqrt{N_{\text{eff}}}}\) * \textbf{Error Estimate}: \(N_{\text{eff }} \approxeq \frac{N}{1 + 2\tau_{Q}}\)

\subsection{Free Energy Calculations}\label{free-energy-calculations}  \textbf{Partition Function} (Free Energy A → B): \(\Delta F_{A \rightarrow B} = - k_{B}T\ln(\frac{Z_{B}}{Z_{A}})\), \(Z = \int d\Gamma e^{- \beta U(\Gamma)}\) 

 \subsubsection{Thermodynamic Free Energy differences}\label{thermodynamic-free-energy-differences}    

\subsubsection{Conformational Free-Energy differences}\label{conformational-free-energy-differences}      

\subsubsection{Alchemical Free energy differences}\label{alchemical-free-energy-differences}          

\subsubsection{Cycles}       \subsection{Systems}\label{systems}  



 \textbf{Hamiltonian \(\equiv\) Newtonian} System with generalized coordinates \(q = \left\{ q_{m} \right\}\), conjugate momenta \(p = \left\{ p_{m} \right\}\), kinetic energy \(K\), potential \(V\). \(H(q,p) = K(q,p) + V(q)\) \textit{Hamilton's equations:} \(\dot{q} = \frac{\partial H}{\partial p},\quad\dot{p} = - \frac{\partial H}{\partial q}\) \textit{Cartesian case:} For \(N\) particles with masses \(m_{n}\), \(H(r,p_{r}) = \sum_{n = 1}^{N}\frac{p_{r,n}^{2}}{2m_{n}} + V(r)\) \({\dot{r}}_{n} = \frac{p_{r,n}}{m_{n}},\quad{\dot{p}}_{r,n} = - \frac{\partial V}{\partial r_{n}} = F_{n}\) \(\Rightarrow\) Hamiltonian \(\Leftrightarrow\) Newton c

\textbf{Code}: \(s(t) = \left\{ N^{-}1\sum_{n = 1}^{N}\left\lbrack {\widetilde{r}}_{n}(t) - r_{n}^{\ast} \right\rbrack^{2} \right\}^{\frac{1}{2}}\)

\raisebox{-5pt}[1em]{\includegraphics{generated/7935136493746420978.pdf}}
\end{document}